\documentclass{article}

\def\npart {IA}
\def\nterm {Michaelmas}
\def\nyear {2023}
\def\nlecturer {Prof N.\ Peake}
\def\ncourse {Vector and Matrices}

\usepackage{alltt}
\usepackage{amsmath, amsfonts, amssymb, amsthm}
\usepackage{booktabs}
\usepackage{caption}
\usepackage{enumitem}
\usepackage{fancyhdr}
\usepackage{graphicx}
\usepackage{mathdots}
\usepackage{mathtools}
\usepackage{microtype}
\usepackage{multirow}
\usepackage{pdflscape}
\usepackage{pgfplots}
\usepackage{siunitx}
\usepackage{textcomp}
\usepackage{slashed}
\usepackage{tabularx}
\usepackage{tikz}
\usepackage{tkz-euclide}
\usepackage[normalem]{ulem}
\usepackage[all]{xy}

\pgfplotsset{compat=1.18}
\pagestyle{fancyplain}

\lhead{\emph{\nouppercase{\leftmark}}}
\rhead{
  \ifnum\thepage=1
  \else
    \npart\ \vline\ \ncourse
  \fi}

\def\nauthor{Marcus Ng}
\author{Based on lectures by \nlecturer \\\small Notes taken by \nauthor}
\date{\nterm\ \nyear}
\title{Part \npart\ --- \ncourse}

\newcommand*{\Cdot}{{\raisebox{-0.25ex}{\scalebox{1.5}{$\cdot$}}}}
\setlist[enumerate,1]{label={(\roman*)}}

\newcommand {\pd}[2][ ]{
  \ifx #1 { }
    \frac{\partial}{\partial #2}
  \else
    \frac{\partial^{#1}}{\partial #2^{#1}}
  \fi
}

% Theorems
\theoremstyle{definition}
\newtheorem*{aim}{Aim}
\newtheorem*{axiom}{Axiom}
\newtheorem*{claim}{Claim}
\newtheorem*{cor}{Corollary}
\newtheorem*{conjecture}{Conjecture}
\newtheorem*{defi}{Definition}
\newtheorem*{eg}{Example}
\newtheorem*{ex}{Exercise}
\newtheorem*{fact}{Fact}
\newtheorem*{law}{Law}
\newtheorem*{lemma}{Lemma}
\newtheorem*{notation}{Notation}
\newtheorem*{prop}{Proposition}
\newtheorem*{question}{Question}
\newtheorem*{problem}{Problem}
\newtheorem*{rrule}{Rule}
\newtheorem*{thm}{Theorem}
\newtheorem*{assumption}{Assumption}
\newtheorem*{assert}{Assertion}

\newtheorem*{remark}{Remark}
\newtheorem*{warning}{Warning}
\newtheorem*{exercise}{Exercise}

\newtheorem{nthm}{Theorem}[section]
\newtheorem{nlemma}[nthm]{Lemma}
\newtheorem{nprop}[nthm]{Proposition}
\newtheorem{ncor}[nthm]{Corollary}

\SetEnumitemKey{cases}{
  label=\underline{Case~\arabic*:},
  ref={\arabic*},
  align=left,
  leftmargin=\parindent
}

% \renewcommand{\labelitemi}{--}
% \renewcommand{\labelitemii}{$\circ$}
% \renewcommand{\labelenumi}{(\roman{*})}

\let\stdsection\section
\renewcommand\section{\newpage\stdsection}

% Strike through
\def\st{\bgroup \ULdepth=-.55ex \ULset}


%%%%%%%%%%%%%%%%%%%%%%%%%
%%%%% Maths Symbols %%%%%
%%%%%%%%%%%%%%%%%%%%%%%%%

% Logical symbols
\newcommand{\contradiction}{
  \begin{tikzpicture}[rotate=45,x=0.5ex,y=0.5ex]
  \draw[line width=.1ex] (0,2) -- (3,2) (0,1) -- (3,1) (1,3) -- (1,0) (2,3) -- (2,0);
  \end{tikzpicture}
}


\let\U\relax
\let\C\relax
\let\G\relax

% Matrix groups
% \newcommand{\GL}{\mathrm{GL}}
% \newcommand{\Or}{\mathrm{O}}
% \newcommand{\PGL}{\mathrm{PGL}}
% \newcommand{\PSL}{\mathrm{PSL}}
% \newcommand{\PSO}{\mathrm{PSO}}
% \newcommand{\PSU}{\mathrm{PSU}}
% \newcommand{\SL}{\mathrm{SL}}
% \newcommand{\SO}{\mathrm{SO}}
% \newcommand{\Spin}{\mathrm{Spin}}
% \newcommand{\Sp}{\mathrm{Sp}}
% \newcommand{\SU}{\mathrm{SU}}
\newcommand{\U}{\mathrm{U}}
% \newcommand{\Mat}{\mathrm{Mat}}

% Matrix algebras
% \newcommand{\gl}{\mathfrak{gl}}
% \newcommand{\ort}{\mathfrak{o}}
% \newcommand{\so}{\mathfrak{so}}
% \newcommand{\su}{\mathfrak{su}}
% \newcommand{\uu}{\mathfrak{u}}
% \renewcommand{\sl}{\mathfrak{sl}}

% Special sets
\newcommand{\C}{\mathbb{C}}
\newcommand{\CP}{\mathbb{CP}}
\newcommand{\GG}{\mathbb{G}}
\newcommand{\N}{\mathbb{N}}
\newcommand{\Q}{\mathbb{Q}}
\newcommand{\R}{\mathbb{R}}
\newcommand{\RP}{\mathbb{RP}}
\newcommand{\T}{\mathbb{T}}
\newcommand{\Z}{\mathbb{Z}}
\renewcommand{\H}{\mathbb{H}}

% Brackets
\newcommand{\abs}[1]{\left\lvert #1\right\rvert}
\newcommand{\bket}[1]{\left\lvert #1\right\rangle}
\newcommand{\brak}[1]{\left\langle #1 \right\rvert}
\newcommand{\braket}[2]{\left\langle #1\middle\vert #2 \right\rangle}
\newcommand{\bra}{\langle}
\newcommand{\ket}{\rangle}
\newcommand{\norm}[1]{\left\lVert #1\right\rVert}
\newcommand{\normalorder}[1]{\mathop{:}\nolimits\!#1\!\mathop{:}\nolimits}
\newcommand{\tv}[1]{|#1|}
\renewcommand{\vec}[1]{\boldsymbol{\mathbf{#1}}}

% not-math
% \newcommand{\bolds}[1]{{\bfseries #1}}
% \newcommand{\cat}[1]{\mathsf{#1}}
% \newcommand{\ph}{\,\cdot\,}
% \newcommand{\term}[1]{\emph{#1}\index{#1}}
% \newcommand{\phantomeq}{\hphantom{{}={}}}

% Probability
% \DeclareMathOperator{\Bernoulli}{Bernoulli}
% \DeclareMathOperator{\betaD}{beta}
% \DeclareMathOperator{\bias}{bias}
% \DeclareMathOperator{\binomial}{binomial}
% \DeclareMathOperator{\corr}{corr}
% \DeclareMathOperator{\cov}{cov}
% \DeclareMathOperator{\gammaD}{gamma}
% \DeclareMathOperator{\mse}{mse}
% \DeclareMathOperator{\multinomial}{multinomial}
% \DeclareMathOperator{\Poisson}{Poisson}
% \DeclareMathOperator{\var}{var}
% \newcommand{\E}{\mathbb{E}}
% \newcommand{\Prob}{\mathbb{P}}

% Algebra
% \DeclareMathOperator{\adj}{adj}
% \DeclareMathOperator{\Ann}{Ann}
% \DeclareMathOperator{\Aut}{Aut}
% \DeclareMathOperator{\Char}{char}
% \DeclareMathOperator{\disc}{disc}
% \DeclareMathOperator{\dom}{dom}
% \DeclareMathOperator{\fix}{fix}
% \DeclareMathOperator{\Hom}{Hom}
% \DeclareMathOperator{\id}{id}
% \DeclareMathOperator{\image}{image}
% \DeclareMathOperator{\im}{im}
% \DeclareMathOperator{\re}{re}
% \DeclareMathOperator{\tr}{tr}
% \DeclareMathOperator{\Tr}{Tr}
% \newcommand{\Bilin}{\mathrm{Bilin}}
% \newcommand{\Frob}{\mathrm{Frob}}

% Others
% \newcommand\ad{\mathrm{ad}}
% \newcommand\Art{\mathrm{Art}}
% \newcommand{\B}{\mathcal{B}}
% \newcommand{\cU}{\mathcal{U}}
% \newcommand{\Der}{\mathrm{Der}}
% \newcommand{\D}{\mathrm{D}}
% \newcommand{\dR}{\mathrm{dR}}
% \newcommand{\exterior}{\mathchoice{{\textstyle\bigwedge}}{{\bigwedge}}{{\textstyle\wedge}}{{\scriptstyle\wedge}}}
% \newcommand{\F}{\mathbb{F}}
% \newcommand{\G}{\mathcal{G}}
% \newcommand{\Gr}{\mathrm{Gr}}
% \newcommand{\haut}{\mathrm{ht}}
% \newcommand{\Hol}{\mathrm{Hol}}
% \newcommand{\hol}{\mathfrak{hol}}
% \newcommand{\Id}{\mathrm{Id}}
% \newcommand{\lie}[1]{\mathfrak{#1}}
% \newcommand{\op}{\mathrm{op}}
% \newcommand{\Oc}{\mathcal{O}}
% \newcommand{\pr}{\mathrm{pr}}
% \newcommand{\Ps}{\mathcal{P}}
% \newcommand{\pt}{\mathrm{pt}}
% \newcommand{\qeq}{\mathrel{``{=}"}}
% \newcommand{\Rs}{\mathcal{R}}
% \newcommand{\Vect}{\mathrm{Vect}}
% \newcommand{\wsto}{\stackrel{\mathrm{w}^*}{\to}}
% \newcommand{\wt}{\mathrm{wt}}
% \newcommand{\wto}{\stackrel{\mathrm{w}}{\to}}
% \renewcommand{\d}{\mathrm{d}}
% \renewcommand{\P}{\mathbb{P}}
%\renewcommand{\F}{\mathcal{F}}

\let\Im\relax
\let\Re\relax

% \DeclareMathOperator{\area}{area}
% \DeclareMathOperator{\card}{card}
% \DeclareMathOperator{\ccl}{ccl}
% \DeclareMathOperator{\ch}{ch}
% \DeclareMathOperator{\cl}{cl}
% \DeclareMathOperator{\cls}{\overline{\mathrm{span}}}
% \DeclareMathOperator{\coker}{coker}
% \DeclareMathOperator{\conv}{conv}
% \DeclareMathOperator{\cosec}{cosec}
% \DeclareMathOperator{\cosech}{cosech}
% \DeclareMathOperator{\covol}{covol}
% \DeclareMathOperator{\diag}{diag}
% \DeclareMathOperator{\diam}{diam}
% \DeclareMathOperator{\Diff}{Diff}
% \DeclareMathOperator{\End}{End}
% \DeclareMathOperator{\energy}{energy}
% \DeclareMathOperator{\erfc}{erfc}
% \DeclareMathOperator{\erf}{erf}
% \DeclareMathOperator*{\esssup}{ess\,sup}
% \DeclareMathOperator{\ev}{ev}
% \DeclareMathOperator{\Ext}{Ext}
% \DeclareMathOperator{\fst}{fst}
% \DeclareMathOperator{\Fit}{Fit}
% \DeclareMathOperator{\Frac}{Frac}
% \DeclareMathOperator{\Gal}{Gal}
% \DeclareMathOperator{\gr}{gr}
% \DeclareMathOperator{\hcf}{hcf}
\DeclareMathOperator{\Im}{Im}
% \DeclareMathOperator{\Ind}{Ind}
% \DeclareMathOperator{\Int}{Int}
% \DeclareMathOperator{\Isom}{Isom}
% \DeclareMathOperator{\lcm}{lcm}
% \DeclareMathOperator{\length}{length}
% \DeclareMathOperator{\Lie}{Lie}
% \DeclareMathOperator{\like}{like}
% \DeclareMathOperator{\Lk}{Lk}
% \DeclareMathOperator{\Maps}{Maps}
% \DeclareMathOperator{\orb}{orb}
% \DeclareMathOperator{\ord}{ord}
% \DeclareMathOperator{\otp}{otp}
% \DeclareMathOperator{\poly}{poly}
% \DeclareMathOperator{\rank}{rank}
% \DeclareMathOperator{\rel}{rel}
% \DeclareMathOperator{\Rad}{Rad}
\DeclareMathOperator{\Re}{Re}
% \DeclareMathOperator*{\res}{res}
% \DeclareMathOperator{\Res}{Res}
% \DeclareMathOperator{\Ric}{Ric}
% \DeclareMathOperator{\rk}{rk}
% \DeclareMathOperator{\Rees}{Rees}
% \DeclareMathOperator{\Root}{Root}
% \DeclareMathOperator{\sech}{sech}
% \DeclareMathOperator{\sgn}{sgn}
% \DeclareMathOperator{\snd}{snd}
% \DeclareMathOperator{\Spec}{Spec}
% \DeclareMathOperator{\spn}{span}
% \DeclareMathOperator{\stab}{stab}
% \DeclareMathOperator{\St}{St}
% \DeclareMathOperator{\supp}{supp}
% \DeclareMathOperator{\Syl}{Syl}
% \DeclareMathOperator{\Sym}{Sym}
% \DeclareMathOperator{\vol}{vol}
\numberwithin{equation}{section}

\begin{document}
\maketitle{
    \small 
    \noindent\textbf{Complex numbers}\\
    Review of complex numbers, including complex conjugate, inverse, modulus, argument and Argand diagram. Informal treatment of complex logarithm, $n$-th roots and complex powers. de Moivre's theorem.\hspace*{\fill}[2]

    \vspace{10pt}
    \noindent\textbf{Vectors}\\
    Review of elementary algebra of vectors in $\R^3$, including scalar product. Brief discussion of vectors in $\R^n$ and $\C^n$; scalar product and the Cauchy-Schwarz inequality. Concepts of linear span, linear independence, subspaces, basis and dimension.

    \vspace{5pt}
    \noindent Suffix notation: including summation convention, $\delta_{ij}$ and $\varepsilon_{ijk}$. Vector product and triple product: definition and geometrical interpretation. Solution of linear vector equations. Applications of vectors to geometry, including equations of lines, planes and spheres.\hspace*{\fill}[5]

    \vspace{10pt}
    \noindent\textbf{Matrices}\\
    Elementary algebra of $3\times 3$ matrices, including determinants. Extension to $n\times n$ complex matrices. Trace, determinant, non-singular matrices and inverses. Matrices as linear transformations; examples of geometrical actions including rotations, reflections, dilations, shears; kernel and image.\hspace*{\fill}[4]

    \vspace{5pt}
    \noindent Simultaneous linear equations: matrix formulation; existence and uniqueness of solutions, geometric interpretation; Gaussian elimination.\hspace*{\fill}[3]

    \vspace{5pt}
    \noindent Symmetric, anti-symmetric, orthogonal, hermitian and unitary matrices. Decomposition of a general matrix into isotropic, symmetric trace-free and antisymmetric parts.\hspace*{\fill}[1]

    \vspace{10pt}
    \noindent\textbf{Eigenvalues and Eigenvectors}\\
    Eigenvalues and eigenvectors; geometric significance.\hspace*{\fill}[2]

    \vspace{5pt}
    \noindent Proof that eigenvalues of hermitian matrix are real, and that distinct eigenvalues give an orthogonal basis of eigenvectors. The effect of a general change of basis (similarity transformations). Diagonalization of general matrices: sufficient conditions; examples of matrices that cannot be diagonalized. Canonical forms for $2 \times 2$ matrices.\hspace*{\fill}[5]

    \vspace{5pt}
    \noindent Discussion of quadratic forms, including change of basis. Classification of conics, cartesian and polar forms.\hspace*{\fill}[1]

    \vspace{5pt}
    \noindent Rotation matrices and Lorentz transformations as transformation groups.\hspace*{\fill}[1]}

\tableofcontents

%%%%%%%%%%%%%%%%%%%%%%
% Lecture 1: 6/10/2023
%%%%%%%%%%%%%%%%%%%%%%
\section{Complex Numbers}
\subsection{Basic properties}
\begin{defi}[Complex number]
    A \emph{complex number} $z \in \C$ is
    \[
        z = a + ib, \quad a, b \in \R, \quad i = \sqrt{-1}  
    \]
    We define two other functions, $\Re()$ and $\Im()$ such that,
    \[
        \Re(z) = a \quad \Im(z) = b 
    \]
    Clearly, we may note that $\R \subset \C$ in the case where $b = 0$
\end{defi}

We can define algebraic manipulation over the complex numbers,
\begin{align*}
    z_1 \pm z_2 &= a_1 \pm a_2 + i(b_1 \pm b_2) \\
    z_1z_2 &= a_1a_2 - b_1b_2 + i(a_1b_2 + a_2b_1)
\end{align*}
We can define the inverse of a complex number
\begin{equation}
    z^{-1} = \frac{1}{a + ib} = \frac{a - ib}{a^2 + b^2}\label{eq:1}
\end{equation}
this demonstrates that $\C$ is \emph{closed} under addition and multiplication. We also define a complex conjugate of $z = a + ib$ as $\overline{z}$ or $z^*$.
\[
    \overline{z}, z^* = a - ib
\]
The modulus of $z$ as $\abs{z}$
\begin{align*}
    \abs{z} &= \abs{a + ib} = \sqrt{a^2 + b^2} \\
    z\overline{z} &= a^2 + b^2 = \abs{z}^2 \\
    z^{-1} &= \frac{\overline{z}}{\abs{z}^2} \tag{see \eqref{eq:1}}
\end{align*}
We also introduce the Argand Diagram, where the coordinate axes are the Real and Imaginary axes. We can see that $\overline{z}$ is the reflection of $z$ in the $x$-axis.

% draw an argand diagram in tikz
\begin{defi}[Triangle inequality]
    For $z_1, z_2 \in \C$, the triangle inequality states that
    \begin{equation}
        \abs{z_1 + z_2} \leq \abs{z_1} + \abs{z_2}\label{eq:2}
    \end{equation}
\end{defi}

\begin{prop}
    The triangle inequality has another common alternative form,
    \begin{equation}
        \abs{z_1 - z_2} \geq \abs{\abs{z_1} - \abs{z_2}}\label{eq:3}
    \end{equation}
\end{prop}

\begin{proof}
    Write $z_1 = z^{\prime}_1 - z^{\prime}_2$, and $z_2 = z^{\prime}_2$ and substitute them into \eqref{eq:2}. Then,
    \begin{align*}
        &\abs{z^{\prime}_1} \leq \abs{z^{\prime}_1 - z^{\prime}_2} + \abs{z^{\prime}_2} \\
        \Rightarrow& \abs{z^{\prime}_1} - \abs{z^{\prime}_2} \leq\abs{z^{\prime}_1 - z^{\prime}_2}\\
        \Rightarrow& \abs{z^{\prime}_1 - z^{\prime}_2} \geq \abs{z^{\prime}_1} - \abs{z^{\prime}_2} \tag{$\star$}
    \end{align*}
    Notice that our labelling of $z_1$ and $z_2$ was arbitrary, and so we can freely swap the assignment. This gives,
    \[
        \abs{z^{\prime}_1 - z^{\prime}_2} \geq \abs{z^{\prime}_2} - \abs{z^{\prime}_1} \tag{$\star\star$}
    \]
    Since $\abs{z^{\prime}_1 - z^{\prime}_2} \geq (\star) \ \text{and} \ (\star\star)$, hence, \eqref{eq:3}
\end{proof}

\begin{defi}[Polar Representation]
    For a complex number $z = a + ib$ we use the argand diagram to write $z = r(\cos{\theta} + i\sin{\theta})$, where
    \[
        r = \abs{z} \quad \text{and} \quad \theta = \arctan\left(\frac{b}{a}\right) = \arg{(z)}  
    \]
\end{defi}

\begin{remark}
    Note that the mapping $(r, \theta) \rightarrow z$ is unique. However, the reverse $r \rightarrow (r, \theta)$ is \emph{NOT} unique since for each value of $\theta$, $\theta + 2n\pi$ is equally valid.

    In order to make this mapping single valued, we will by convention restrict $\theta$ to the principle value, where
    \[
        -\pi < \theta \leq \pi  
    \]

    This multivalued property becomes quite interesting when we consider complex functions. We can imagine taking the multivalued imaginary $\ln$ function as a multistory carpark, walking round and round up and down. To fix this what we do is to build a brick wall on the upramp (``branch cut'').
\end{remark}

Using complex numbers in this form allows us to view complex arithmetic geometrically. For instance, let
\begin{align*}
    z_k &= x_k + iy_k \quad \text{for} \quad k = 1, 2 \\
    &= r_k(\cos{\theta_k} + i\sin(\theta_k))
\end{align*}
\begin{align*}
    z_1z_2 &= r_1(\cos{\theta_1} + i\sin{\theta_1})r_2(\cos{\theta_2} + i\sin{\theta_2}) \\
    &= r_1r_2[(\cos{\theta_1}\cos{\theta_2} - \sin{\theta_1}\sin{\theta_2}) + i(\cos{\theta_1}\sin{\theta_2} + \sin{\theta_1}\cos{\theta_2})] \\
    &= r_1r_2(\cos{(\theta_1 + \theta_2)} + i\sin{(\theta_1 + \theta_2)}) \numberthis \label{eq:4}
\end{align*}

\subsection{Complex exponential function}
\begin{defi}[Complex exponential function]
    \begin{align*}
        \exp(z) &= 1 + z + \frac{z^2}{2!} + \cdots + \frac{z^n}{n!} + \cdots \\
        &=\sum_{n=0}^{\infty}{\frac{z^n}{n!}} \numberthis \label{eq:5}
    \end{align*}
\end{defi}
This function is ``uniformly convergent'' (something that we will cover next term in analysis). That means that we are able to do lots of algebraic manipulation without too much hesitation. =

\begin{prop}[Multiplication of $\exp()$]
    \[
        \exp(z_1)\exp(z_2) = \exp(z_1 + z_2)  
    \]
\end{prop}
\begin{proof}
    We will consider the double sum
    \begin{align*}
        \sum_{m,n=0}^{\infty}{a_mn} &= a_{00} + a_{01} + a_{02} + a_{03} + \cdots \\
         &+ a_{10} + a_{11} + a_{12} + \cdots \\
         &+ a_{20} + a_{21} + a_{22} + \cdots \\
         &+ a_{30} + a_{31} + a_{32} + \cdots \\
         &\vdots
    \end{align*}
    Then by collecting terms along the diagonal we can regroup the sum into
    \[
        =\sum_{r=0}^{\infty}\sum_{m=0}^{r}{a_{r-m, m}}
    \]
    Hence, 
    \begin{align*}
        \exp(z_1) \cdot \exp(z_2) &= \sum_{r=0}^{\infty}\sum_{m=0}^{r}{\frac{z_1^{r-m}}{(r-m)!}\frac{z_2^{m}}{m!}} \\
        &= \sum_{r=0}^{\infty}{\frac{1}{r!}}\sum_{m=0}^{r}{\frac{r!}{(r-m)!m!}z_1^{r-m}z_2^{m}} \\
        &= \sum_{r=0}^{\infty}{\frac{1}{r!} (z_1 + z_2)^r} \\
        &= \exp(z_1 + z_2)
    \end{align*}
\end{proof}

%%%%%%%%%%%%%%%%%%%%%%
% Lecture 2: 8/10/2023
%%%%%%%%%%%%%%%%%%%%%%
Similarly, we can define the complex $\sin(z)$ and $\cos(z)$ functions.
\begin{defi}[Complex sine and cosine]
    \begin{align*}
        \sin{z} &= \sum_{n=0}^{\infty}{(-1)^n\frac{z^{2n+1}}{(2n+1)!}} \\
        \cos{z} &= \sum_{n=0}^{\infty}{(-1)^n\frac{z^{2n}}{(2n)!}} \numberthis \label{eq:6}
    \end{align*}
\end{defi}

From \eqref{eq:5} we can derive that
\begin{align*}
    e^{iz} &= \sum_{n=0}^{\infty}{\frac{i^nz^n}{n!}} \\
    &= \sum_{n=0}^{\infty}{(-1)^n\frac{z^{2n+1}i^{2n+1}}{(2n+1)!}} + \sum_{n=0}^{\infty}{(-1)^n\frac{z^{2n}i^{2n}}{(2n)!}} \\
    &= i\sin{z} + \cos{z} \numberthis \label{eq:7}
\end{align*}

We may consider $\theta \in \R$ which gives,
\begin{equation}
    z = re^{i\theta} = r(\cos{\theta} + i\sin{\theta}) \label{eq:8}
\end{equation}

\subsection{Roots of unity}
We seek to solve equations of the form $z^n = 1$ where $n \in \Z$. Since this is a polynomial of degree $n$ we expect to find $n$ solutions. By applying the exponential form we found earlier,
\begin{align*}
    z^n = 1 &= \exp(i2 \pi k), \quad k \in \Z \\
    z &= \exp(i2 \pi \frac{k}{n})
\end{align*}
Notice that the roots will repeat after $k = n-1$. Hence the solutions are
\[
    z = \exp(i2 \pi \frac{k}{n}), \quad k = 0, 1, 2, \cdots, n-1
\]

\begin{eg}[$ z^7 = 2$]
    \[
        z = 2^{\frac{1}{7}} \exp(i 2\pi \frac{m}{7}), \quad m = 0, 1, 2, \cdots, 6
    \]
\end{eg}

An interesting point to note is that for the polynomial $z^n - 1 = 0$, if we arbitrarily choose a root $w = \exp(\frac{2\pi i}{n})$, then all the remaining roots will be
\[
    1, w, w^2, w^3, \cdots, w^{n-1}
\]
Further, if we note that since the coefficient of $z^{n-1}$ is $0$, that implies that the sum of the roots must be $0$.
\[
   1 + w + w^2 + w^3 + \cdots + w^{n-1} = 0
\]

\subsection{Complex Logarithm and Powers}
\begin{defi}[Complex Log]
    The \emph{complex log}, $\log{z}$, is a solution of the following
    \[
        e^w = z  
    \]
    Hence, $\exp()$ and $\log()$ are inverse functions, i.e;
    \[
        e^{\log{z}} = z
    \]
    If we write $z = re^{i\theta}$ from \eqref{eq:8} then,
    \begin{align*}
        \log{z} = \log{re^{i\theta}} &= \log{r} + \log{e^{i\theta}} \\
        &= log{\abs{z}} + i\arg{(z)} \numberthis \label{eq:9}
    \end{align*}
    It is important to note that this is not a single valued function. In order to define a single valued $\log$ we will insist that
    \[
        -\pi < \arg{(z)} \leq \pi
    \]
    Hence,
    \begin{equation}
        \log{z} = log{\abs{z}} + i\theta, \quad -\pi < \theta \leq \pi \label{eq:10}
    \end{equation}
\end{defi}

\begin{eg}
    \begin{align*}
        \log{(2i)} &= \log{2e^{i\frac{\pi}{2}}} \\
        &=\log{2} + i\frac{\pi}{2}
    \end{align*}
\end{eg}

\begin{defi}[Complex Power]
    Since we have defined the complex logarithm we can now define the complex power, $z^\alpha$ for $z, \alpha \in \C$ as
    \begin{equation}
        z^\alpha = e^{\alpha \log{z}} \label{eq:11}
    \end{equation}
    Again, this function is multivalued. Specifically using \eqref{eq:10},
    \[
        z^\alpha = e^{\alpha \log{z}}e^{i \alpha \theta}e^{i \alpha 2\pi n}, \quad n \in \Z
    \]
    However, unlike the roots of unity, there is no reason why this series will repeat and terminate. If $\alpha$ is real and rational then will there be finite solutions. Once again, we will make this function single valued by insisting on using the principle value for $\theta$
    \[
        -\pi < \arg{(z)} \leq \pi
    \]
\end{defi}

\begin{ex}\leavevmode
    \begin{enumerate}
        \item \begin{align*}
            (1 + i)^\frac{1}{2} &= e^{\frac{1}{2}\sqrt{2}}e^\frac{i\pi}{8}e^{in\pi}, \quad n = 0, 1 \\
            &= 2^\frac{1}{4}e^\frac{i\pi}{8}, 2^\frac{1}{4}e^\frac{i9\pi}{8} \\
            &= 2^\frac{1}{4}e^\frac{i\pi}{8} \tag{by insiting on the principle value}
        \end{align*}
        \item \begin{align*}
            i^i &= \exp(i \log{(e^\frac{i\pi}{2})}) \tag{from \eqref{eq:11}} \\
            &= \exp(i \frac{i\pi}{2}) \\
            &= \exp(-\frac{\pi}{2})
        \end{align*}
    \end{enumerate}
\end{ex}

\subsection{De Moivre's Theorem}
\begin{thm}[De Moivre's Theorem]
    \begin{equation}
        \cos{n\theta} + i\sin{n\theta} = (\cos{\theta} + i\sin{\theta})^n, \quad n \in \Z \label{eq:12}
    \end{equation}
\end{thm}

\begin{proof}
    By induction, we will first consider the case where $n > 0$. Clearly, $n=1$ is true. Assume $n=k$ is true for $k > 1$.
    \begin{align*}
        (\cos{\theta} + i\sin{\theta})^{k+1} &= (\cos{k\theta} + i\sin{k\theta})(\cos{\theta} + i\sin{\theta}) \\
        &= \cos{k\theta}\cos{\theta} - \sin{k\theta}\sin{\theta} + i(\cos{k\theta}\sin{\theta} + \sin{k\theta}\cos{\theta}) \\
        &= \cos{(k+1)\theta} + i\sin{(k+1)\theta}
    \end{align*}
    Now when $n < 0$, we write $m = -n$ such that $m > 0$
    \begin{align*}
        (\cos{\theta} + i\sin{\theta})^{-m} &= (\cos{m\theta} + i\sin{m\theta})^{-1} \\
        &= \cos{m\theta} - i\sin{m\theta} \tag{by ($\star$)}\\
        &= \cos{(-m\theta)} - i\sin{(-m\theta)}
    \end{align*}
    Note $(\star)$:
    \[
        z^{-1} = \frac{z^*}{\abs{z}^2}    
    \]
\end{proof}

Using De Moivre's theorem we can work out multiple angle formulae by comparing real and imaginary parts:
\begin{eg}
    \begin{align*}
        \cos{5 \theta} + i\sin{5 \theta} &= (\cos{\theta} + i\sin{\theta})^5 \\
        &\vdots \\
        \cos{5 \theta} &= 5\cos{\theta} - 20 \cos^3{\theta} + 16\cos^5{\theta}\\
        \sin{5 \theta} &= 5\cos{\theta} - 20 \cos^3{\theta} + 16\cos^5{\theta}\\
    \end{align*}
\end{eg}
Note that for non-integer powers of $n$ we will have to insist on the principle value \eqref{eq:12}

\subsection{Lines and Circles in $\C$}
\begin{defi}[Straight Line]
    A line $l$ through $z_0 \in \C$ parallel to $w \in \C$ can be generally written
    \[
        z = z_0 + \lambda w, \quad \lambda \in \R    
    \]
    However, it would be better to get rid of the parameterisation. We can use the fact that $\lambda$ must be real and hence its conjugate is equal to itself.
    \begin{align*}
        \lambda &= \frac{z - z_0}{w} \in \R\\
        &\Rightarrow \frac{z - z_0}{w} = \frac{\overline{z} - \overline{z_0}}{\overline{w}} \\
        &\Leftrightarrow z \overline{w} - \overline{z}w = z_0 \overline{w} - \overline{z_0} w
    \end{align*}
\end{defi}

\begin{defi}[Circle]
    A circle can be defined from its center $c \in \C$ and radius $\rho \in \R$. By definition, if $z$ lies on the circle,
    \begin{align*}
        &\abs{z - c} = \rho \\
        \Rightarrow& (z - c)(\overline{z} - \overline{c}) = \rho^2 \\
        \Leftrightarrow& z\overline{z} - \overline{c}z - c\overline{z} = \rho^2 - c\overline{c}
    \end{align*}
\end{defi}

%%%%%%%%%%%%%%%%%%%%%%
% Lecture 3: 11/10/2023
%%%%%%%%%%%%%%%%%%%%%%

\section{Vectors}
\subsection{Definition \& basic properties}
\begin{defi}[Vector]
    A vector $\vec{x}$ has a positive length and a direction. If,
    \[
        \abs{\vec{x}} = 0 \Rightarrow \vec{x} = \vec{0}  
    \]
    If $\abs{\vec{x}} = 1$ then we write $\uvec{x}$ as \emph{unit vector}. Note the distinction between scalar $0$ and vector $\vec{0}$.
\end{defi}

\begin{defi}[Vector Field]
    A vector field $F = F(\vec{x})$ is a vector function of position and maybe time.
\end{defi}

\begin{defi}[Vector addition]
    % insert paralellogram diagram.
    We can define the addition of vectors using the paralellogram rule,
    \begin{align*}
        \vec{a} + \vec{b} &= \vec{c} \\
        \pvec{OA} + \pvec{OB} &= \pvec{OC}
    \end{align*}
    More abstractly, addition will have the following properties
    \begin{enumerate}[eqn]
        \item $\vec{a} + \vec{b} = \vec{b} + \vec{a}$ \hspace*{\fill}(commutativity) \label{eq:2-1}
        \item $\vec{a} + (\vec{b} + \vec{c}) = (\vec{a} + \vec{b}) + \vec{c}$ \hspace*{\fill}(associativity) \label{eq:2-2}
        \item $\exists \vec{0}, \quad \vec{a} + \vec{0} = \vec{a}$ \hspace*{\fill}(unique zero) \label{eq:2-3}
        \item $\vec{a} + -\vec{a} = 0$ \hspace*{\fill}(unique additive inverse) \label{eq:2-4}
    \end{enumerate}
\end{defi}
\begin{defi}[Multiplication by a scalar]
    For $\lambda \in \R$ (we will do complex scalars later in the course), $\lambda\vec{a}$ has magnitutde $\abs{\lambda}\abs{\vec{a}}$ and is parallel/anti-parallel to $\vec{a}$ if $\lambda<0$ \textbackslash $\lambda>0$. 
    Likewise, we have the following properties
    \begin{enumerate}[eqn]
        \item $(\lambda + \mu)\vec{a} = \lambda\vec{a} + \mu\vec{a}$ \label{eq:2-5}
        \item $\lambda(\vec{a} + \vec{b}) = \lambda\vec{a} + \lambda\vec{b}$ \label{eq:2-6}
        \item $\lambda(\mu\vec{a}) = (\lambda\mu)\vec{a}$ \label{eq:2-7}
        \item $1\vec{a} = \vec{a}$ \label{eq:2-8}
    \end{enumerate}
\end{defi}

Implicit in the two definitions is that the set of vectors is closed under addition and scalar multiplication.

\begin{defi}[Vector Space]
    $V$ is a \emph{vector space} if it is a set of vectors which satisfy properties \eqref{eq:2-1} -- \eqref{eq:2-8} inclusive and is closed under addition and multiplication. Note that for now, we will be considering the vector space $V$ over the real numbers $\R$, but we will consider complex fields later. 
\end{defi}

\begin{eg}[$\R^n$ is a vector space]
\end{eg}
\begin{proof}
    \[
        \R^n = \{\vec{x} = (x_1, x_2, \ldots, x_n) \mid x_j \in \R, \ 1 \leq j \leq n\}
    \]
    Consider $\vec{x} , \vec{y} \in \R$
    \[
        \vec{x} + \vec{y} = (x_1 + y_1, x_2 + y_2, \ldots, x_n + y_n,)
    \]
    From this we see that property \eqref{eq:2-1} and \eqref{eq:2-2} is satisfied as we can inherit them from the Real numbers.

    \begin{enumerate}
        \item $\vec{0} = (0, 0, \ldots, 0)$ \hspace*{\fill}\eqref{eq:2-3}
        \item $-\vec{x} = (-x_1, -x_2, \ldots, -x_n)$ \hspace*{\fill}\eqref{eq:2-4}
        \item $\lambda\vec{x} = (\lambda x_1, \lambda x_2, \ldots, \lambda x_n)$ \hspace*{\fill}\eqref{eq:2-5} - \eqref{eq:2-8}
    \end{enumerate}
\end{proof}

Geometrically, we may interprit $\R^1$ as a line, and $\R^2$ as a plane. However, notice that the converse is not always true. Due to \eqref{eq:2-3} the space has to contain $0$. Therefore, if the line or plane does not go through the origin, it is not a vector space.

\subsection{Scalar Product}
\begin{defi}[Geometric definition in $\R^2$, $\R^3$]
    % insert diagram of dot product
    Given two vectors $\vec{a}, \vec{b}$ we define
    \[
        \vec{a} \cdot \vec{b} = \abs{\vec{a}}\abs{\vec{b}}\cos{\theta}  
    \]
    where $\theta$ is the angle between the two vectors.
\end{defi}
From this, we can see that
\begin{enumerate}
    \item $\vec{a} \cdot \vec{b} = \vec{b} \cdot \vec{a}$ 
    \item $\vec{a} \cdot \vec{a} = 0 \Leftrightarrow \vec{a} = \vec{0}$
    \item $\vec{a} \cdot \vec{b} = 0 \Leftrightarrow$ $\vec{a}$ and $\vec{b}$ are orthogonal
    \[\numberthis\]
\end{enumerate}

\begin{defi}[Projection]
    % insert diagram of projection
    $\vec{b}^\perp$ is the projection of $\vec{b}$ onto $\vec{a}$. Specifically, 
    \begin{equation}
        \vec{b}^\perp = \abs{\vec{b}}\cos{\theta}\frac{\vec{a}}{\abs{\vec{a}}} = \frac{(\vec{b}\cdot\vec{a})\vec{a}}{(\vec{a}\cdot\vec{a})} = (\uvec{a} \cdot \vec{b})\uvec{a} 
    \end{equation}
\end{defi}

We will move to a most abstract defintion of a scalar product
\begin{defi}[Inner product]
    We define $\vec{x} \cdot \vec{y} = \braket{\vec{x}}{\vec{y}} \in \R$ as an \emph{inner product}, if it satisfies
    \begin{enumerate}[eqn]
        \item $\vec{x} \cdot \vec{y} = \vec{y} \cdot \vec{x}$ \hspace*{\fill}(Symmetry) \label{eq:2-11}
        \item $\vec{x}, \vec{y}, \vec{z} \in V, \quad, \lambda, \mu \in \R$ \hspace*{\fill}(Linearity in the second argument) \label{eq:2-12}
        \[
            \vec{x} \cdot (\lambda\vec{y} + \mu\vec{z}) = \lambda\vec{x}\vec{y} + \mu\vec{x}\vec{z}  
        \]
        \item $\vec{x} \cdot \vec{x} \geq 0 $ and $\vec{x} \cdot \vec{x} = 0 \Leftrightarrow \vec{x} = \vec{0}$ \hspace*{\fill}(Positive definite) \label{eq:2-13}
    \end{enumerate}

    Note that we can use \eqref{eq:2-11} together with \eqref{eq:2-12} to derive linearlity in the first argument,
    \begin{equation}
        (\lambda\vec{y} + \mu\vec{z}) \cdot \vec{x} = \lambda\vec{x}\vec{y} + \mu\vec{x}\vec{z} \label{eq:2-15}
    \end{equation}
    We will see later in the course that linearity in the first argument does not hold for vector spaces over the complex field.

    We can also define the \emph{norm} of a vector $\vec{a}$ as 
    \[
        (\vec{a} \cdot \vec{a})^\frac{1}{2} = \abs{\vec{a}} = \norm{\vec{a}}  
    \]
\end{defi}

\begin{remark}
    Using this abstracted definition of an inner product, we can define an inner product between functions. For example define for function $f$ and $g$
    \[
        \braket{f}{g} = \int_{0}^{1}{fg \ \diffd x}
    \]
    This definition satisfies \eqref{eq:2-11} and \eqref{eq:2-12} and we can check \eqref{eq:2-13}
    \[
        \braket{f}{f} = \int_{0}^{1}{f^2 \ \diffd x} \geq 0
    \]
    and further if 
    \[
        \int_{0}^{1}{f^2 \ \diffd x} = 0 \Rightarrow f = 0
    \]
\end{remark}

\subsection{Cauchy-Schwartz Inequality}
\begin{thm}[CS Inequality]
    Given two vectors $\vec{x}$, $\vec{y}$
    \begin{equation}
        \vec{x} \cdot \vec{y} \leq \abs{\vec{x}}\abs{\vec{y}} \quad \vec{x}, \vec{y} \in \R^n
    \end{equation}
    with equality if and only if $\vec{x} = 0$ or $\vec{y} = 0$ or $\vec{x} = \lambda \vec{y}$ for some $\lambda \in \R$.
\end{thm}
\begin{proof}
    Consider $\abs{\vec{x} - \lambda \vec{y}}^2 \geq 0$
    \begin{align*}
        &(\vec{x} - \lambda \vec{y}) \cdot (\vec{x} - \lambda \vec{y}) \geq 0 \\
        \Rightarrow& \lambda^2 \vec{y} \cdot \vec{y} - 2\lambda\vec{x}\vec{y} + \vec{x} \cdot \vec{x} \geq 0
    \end{align*}
    Suppose $\vec{x}, \vec{y} \neq \vec{0}$, we have a quadratic equation in $\lambda$. Therefore since the quadratic is greater than or equal to $0$, the discriminant must be less than or equal to $0$.
    \begin{align*}
        \Delta &= [2(\vec{x} \cdot \vec{y})]^2 - 4 (\vec{y} \cdot \vec{y})(\vec{x} \cdot \vec{x}) \leq 0 \\
        \Rightarrow& [2(\vec{x} \cdot \vec{y})]^2 \leq 4 (\vec{y} \cdot \vec{y})(\vec{x} \cdot \vec{x}) \\
        \Rightarrow& (\vec{x} \cdot \vec{y})^2 \leq (\vec{y} \cdot \vec{y})(\vec{x} \cdot \vec{x}) \\
        \Rightarrow& (\vec{x} \cdot \vec{y})^2 \leq \abs{\vec{x}}^2\abs{\vec{y}}^2
    \end{align*}
    Furthermore, we have equality when
    \[
        \abs{\vec{x} - \lambda \vec{y}}^2 = 0 \Leftrightarrow \vec{x} = \lambda \vec{y}
    \]
    Hence the result
\end{proof}

\begin{eg} $\vec{x} = (\alpha, \beta, \gamma), y = (1, 1, 1)$
    \begin{align*}
        \alpha + \beta + \gamma &\leq \sqrt{\alpha^2 + \beta^2 + \gamma^2} \sqrt{3} \tag{\ref{eq:2-15}} \\
        \alpha^2 + \beta^2 + \gamma^2 &\geq \alpha\beta + \beta\gamma + \gamma\alpha \tag{rearranging}
    \end{align*}
\end{eg}

Note that we must not use the geometric definition of the dot product to prove the CS inequality because $\cos{\theta}$ is not defined for higher dimensions. However, we can use CS to define an angle in $\R^n$
\[
    \theta = \arccos{\left(\frac{\vec{x}\vec{y}}{\abs{\vec{x}}\abs{\vec{y}}}\right)}    
\]
since CS gives us that the fraction is less than or equal to $1$.

%%%%%%%%%%%%%%%%%%%%%%
% Lecture 4: 13/10/2023
%%%%%%%%%%%%%%%%%%%%%%

A further consequence of \eqref{eq:2-15} is the triangle inequality \eqref{eq:2}. We will consider the following
\begin{align*}
    \abs{\vec{x} + \vec{y}}^2 &= (\vec{x} + \vec{y}) \cdot (\vec{x} + \vec{y}) \\
    &= \vec{x} \cdot \vec{x} + 2 \vec{y} \cdot \vec{x} + \vec{y} \cdot \vec{y} \tag{collect terms over the reals} \\
    &= \abs{\vec{x}}^2 + \abs{\vec{y}}^2 + 2 \vec{y} \cdot \vec{x} \\
    &\leq \abs{\vec{x}}^2 + \abs{\vec{y}}^2 + 2 \abs{\vec{y} \cdot \vec{x}} \\
    &\leq \abs{\vec{x}}^2 + \abs{\vec{y}^2 + 2\abs{\vec{x}}}\abs{\vec{y}} \tag{By CS}\\
    &= (\abs{\vec{x}} + \abs{\vec{y}})^2 
\end{align*}
We may then take the positive square root on both sides to arive at
\begin{equation}
    \abs{\vec{x} + \vec{y}} \leq \abs{\vec{x}} + \abs{\vec{y}}
\end{equation}

\subsection{Vector Product}
We will see later in the course that this extends to the wedge product in higher dimensions. For now, we will define the product in $\R^3$.
\begin{defi}[Vector Product]
    Consider $\vec{a}, \vec{b} \in \R^3$. Define the vector product
    \[
        \vec{a} \times \vec{b} = \abs{\vec{a}}\abs{\vec{b}}\sin{\theta}\uvec{n}
    \]
    Where the direction of $\uvec{n}$ is such that $\vec{a}$, $\vec{b}$ and $\uvec{n}$ for a \emph{right handed system}. That is, when I curl my right hand fingers in the direction of $\theta$, the thumb points in the direction of $\uvec{n}$.
\end{defi}

Here are some properties of the vector product
\begin{enumerate}
    \item $\vec{a} \times \vec{b} - = \vec{b} \times \vec{a}$ \hspace*{\fill}(Anti-symmetric)
    \item $\vec{a} \times \vec{a} = \vec{0}$
    \item If $\vec{a} \times \vec{b} = \vec{0} \Rightarrow \vec{a} = \mu\vec{b}, \quad \text{for} \mu \in \R$
    \item $\vec{a} \times (\lambda \vec{b}) = \lambda (\vec{a} \times \vec{b})$
    \item[] \hspace*{\fill}\ilnumberthis
\end{enumerate}

\begin{defi}[Vector Area]
    We define the \emph{vector area} between two vectors $a$ and $b$ to be
    \[  
        \text{Vector Area} \equiv \frac{1}{2}(\vec{a} \times \vec{b})
    \]

    We may consider as an example the regular area of  a triangle, which can be found
    \begin{align*}
        \text{A} &= \frac{1}{2}\abs{\pvec{OB}}\abs{\pvec{BN}} \\
        &=  \frac{1}{2}\abs{\pvec{OB}}\abs{\pvec{OB}}\sin{\theta} \\
        &= \frac{1}{2}\abs{\vec{a} \times \vec{b}}
    \end{align*}
    
    Note that to turn this into a vector area, we omit the absolute value and use the vector itself to represent the area. An interesting consequence of this is that although the surface area of a sphere is $4\pi r^2$ the it's vector area is $\vec{0}$!
\end{defi}

\subsection{Scalar Triple Product}

\begin{defi}[Scalar Triple Product]
    Given three vectors $\vec{a}, \vec{b}, \vec{c} \in \R^3$, the triple product is defined as
    \begin{equation}
        \vec{a} \cdot (\vec{b} \times \vec{c}) = (\vec{b} \times \vec{c}) \cdot a = -a \cdot (c \times b)
    \end{equation}
    clearly, this is a scalar. Geometrically, we can interprit the triple product as the volume of a parallelopiped formed by $a$, $b$ and $c$, where the three vectors form a right handed system.
    
    % insert diagram of parallelopiped
    \begin{align*}
        \text{Volume} \ &= \ \text{Base Area} \ \times \perp \text{height} \\
        &\Rightarrow \abs{\vec{b} \times \vec{c}} \abs{\vec{a}} \sin{\theta} \\
        &= \abs{\vec{b} \times \vec{a}}\abs{\vec{a}}\cos{\bkt{\frac{\pi}{2} - \theta}} \\
        &= \abs{(\vec{b} \times \vec{a}) \cdot \vec{a}}
    \end{align*}
    Furthermore, since we specified that the vectors form a right handed system,
    \[
        \vec{a} \cdot (\vec{b} \times \vec{c}) \geq 0
    \]
\end{defi}

\begin{notation}
    We may instead write the tripple product using square brackets as follows
    \begin{equation}
        [\vec{a}, \vec{b}, \vec{c}] = [\vec{b}, \vec{c}, \vec{a}] =  [\vec{c}, \vec{a}, \vec{b}]
    \end{equation}

    where all even permutations are the equal, (even permutation meaning permutations performing an even number of pariwise swaps)
    \begin{equation}
        [\vec{a}, \vec{b}, \vec{c}] = - [\vec{b}, \vec{a}, \vec{c}] = - [\vec{a}, \vec{c}, \vec{b}] = - [\vec{c}, \vec{b}, \vec{a}]
    \end{equation}
    and all odd permutations are negative of the triple product.
\end{notation}

\begin{remark}
    The triple product can tell us if the three vectors are coplanar. Specifically,
    \begin{equation}
        \vec{a}, \vec{b}, \vec{c} \ \text{coplanar} \Leftrightarrow [\vec{a}, \vec{b}, \vec{c}] = 0  
    \end{equation}
\end{remark}

\begin{cor}
    Any triple product where two of the vectors are the same is equal to $0$.
    \[
        [\vec{a}, \vec{a}, \vec{c}] = [\vec{a}, \vec{b}, \vec{b}] = \ldots = 0  
    \]
\end{cor}

\subsection{Spanning sets \& Basis}
We will build up from the smaller example of $\R^2$, to $\R^3$, before working in $\R^n$. In general, spanning sets together with linear independance gives us a basis.

\subsubsection*{Working in $\R^2$}
\begin{defi}[Spanning sets]
    Consider $2$ vectors $\vec{a}, \vec{b} \neq 0 \in \R^2$ with $\vec{a} \times \vec{b} \neq 0$ i.e; $\vec{a}$ and $\vec{b}$ are not parallel. Then, any other vector $\vec{r}$ in this plane can be written as 
    \begin{equation}
        \vec{r} = \lambda \vec{a} + \mu \vec{b}, \quad \text{for} \lambda, \mu \in \R \label{eq:2-22}
    \end{equation}
    Then we say that the set $\{\vec{a}, \vec{b}\}$ spans $\R^2$, and every other vector in the space can be written as a linear combination of $\vec{a}$ and $\vec{b}$.
\end{defi}

\begin{prop}
    \eqref{eq:2-22} produces a unique representation.
\end{prop}

\begin{proof}
    Suppose 
    \[
        r = \lambda \vec{a} + \mu \vec{b} = \lambda' \vec{a} + \mu' \vec{b}    
    \]
    a useful trick to employ it to apply cross product with one of the terms, allowing us to remove it. Taking $\times \vec{b}$ to both sides of the equation we get
    \[
        (\lambda - \lambda')\vec{a} \times \vec{b} = \vec{0}  
    \]
    finally, as we have defined $\vec{a}$ and $\vec{b}$ to be non parallel,
    \[
        \therefore \lambda - \lambda' = 0
    \]

    Similarly, we can multiply $\times \vec{a}$ to arrive at
    \[
        \mu - \mu' = 0  
    \]
    Hence \eqref{eq:2-22} is unique.
\end{proof}

\begin{defi}[Linear Independance]
    Two vectors $\vec{a}, \vec{b}$ are said to be \emph{linearly independant} if, $\forall \alpha, \beta \in \R$
    \begin{equation}
        \alpha \vec{a} + \beta \vec{b} = 0 \Rightarrow \alpha = \beta = 0 \label{eq:2-23}
    \end{equation}
    Later on, we will see that this is generalised into $\R^n$
\end{defi}

\begin{remark}
    $\vec{a}$ and $\vec{b}$ are linearly dependant in $\R^2$ if $\vec{a} \times \vec{b} \neq 0$. To see that it is the case, take
    \begin{align*}
        \vec{b} \times \eqref{eq:2-23} &\Rightarrow \alpha (\vec{b} \times \vec{a}) = 0 \\
        \vec{a} \times \eqref{eq:2-23} &\Rightarrow \beta (\vec{a} \times \vec{b}) = 0 \\
    \end{align*}
    In both cases, since $\vec{a} \times \vec{b} \neq 0$, $\alpha = \beta = 0$
\end{remark}

\begin{defi}[Basis]
    A set $\{\vec{a}, \vec{b}\}$ is a basis of $\R^n$ if it spans $\R^2$ and it is linearly independant.
\end{defi}

\begin{remark}
    Note that the number of dimensions of a space in fact follows from the size of the set of minimal basis vectors.
\end{remark}

\begin{eg}
    \begin{enumerate}
        \item $\{\uvec{i}, \uvec{j}\}$ is a basis of $\R^2$
        \item $\{\uvec{i}, \uvec{j}, \vec{e}\}$ where $\vec{e} = \uvec{i} + \uvec{j}$ certainly spans $\R^2$, but it does not form a basis since the three vector system is not linearly independant. To see take,
        \[
            \uvec{i} + \uvec{j} + \vec{e} = \vec{0}    
        \]
    \end{enumerate}
\end{eg}

\begin{prop}
    A basis is the smallest spanning set
\end{prop}

\begin{proof}
    See Example sheet exercise
\end{proof}

\subsubsection*{Working in $\R^3$}
\begin{defi}[Basis in $\R^3$]
    If $\vec{a}, \vec{b}, \vec{c}$ are non-coplanar, then $\{\vec{a}, \vec{b}, \vec{c}\}$ is a basis of $\R^3$. We can write any vector $r$ as 
    \begin{equation}
        \vec{r} = \lambda \vec{a} + \mu \vec{b} + \nu \vec{c}, \quad \text{for}, \lambda, \mu, \nu \in \R \label{eq:2-24}
    \end{equation}
\end{defi}
\begin{proof}
    To show that it is uniquely defined, we will proceed by construction of the variables. Using a similar trick, using the tripple product to get rid of terms as repeated vectors in the triple product go to zero. $\eqref{eq:2-24} \cdot (\vec{b} \times \vec{c})$
    \begin{align*}
        \vec{r} \cdot (\vec{b} \times \vec{c}) &= \lambda \vec{a} \cdot (\vec{b} \times \vec{c}) \\
        &\Rightarrow \lambda = \frac{\vec{r} \cdot (\vec{b} \times \vec{c})}{\vec{a} \cdot (\vec{b} \times \vec{c})} = \frac{\vec{r} \cdot (\vec{b} \times \vec{c})}{[\vec{a}, \vec{b}, \vec{c}]}\\
        \vec{r} \cdot (\vec{c} \times \vec{a}) &= \mu \vec{b} \cdot (\vec{c} \times \vec{a}) \\
        &\Rightarrow \mu = \frac{\vec{r} \cdot (\vec{c} \times \vec{a})}{\vec{b} \cdot (\vec{c} \times \vec{a})}  = \frac{\vec{r} \cdot (\vec{c} \times \vec{a})}{[\vec{a}, \vec{b}, \vec{c}]} \\
        \vec{r} \cdot (\vec{a} \times \vec{b}) &= \nu \vec{c} \cdot (\vec{a} \times \vec{b}) \\
        &\Rightarrow \nu = \frac{\vec{r} \cdot (\vec{a} \times \vec{b})}{\vec{c} \cdot (\vec{a} \times \vec{b})}  = \frac{\vec{r} \cdot (\vec{a} \times \vec{b})}{[\vec{a}, \vec{b}, \vec{c}]} \\
    \end{align*}
    Hence, $\lambda, \mu, \nu$ can be determined for any $\vec{r}$ and $\{\vec{a}, \vec{b}, \vec{c}\}$ spans $\R^3$.
%%%%%%%%%%%%%%%%%%%%%%
% Lecture 5: 16/10/2023
%%%%%%%%%%%%%%%%%%%%%%
    Further to show linear independance, let 
    \[
        \vec{r} = \alpha \vec{a} + \beta \vec{b} + \gamma \vec{c} \quad \text{for} \ [\vec{a}, \vec{b}, \vec{c}] \neq 0
    \]
    We apply the same trick, using the triple product
    \begin{align*}
        \vec{r} \cdot (\vec{b} \times \vec{c}) &\Rightarrow \alpha [\vec{a}, \vec{b}, \vec{c}] = 0 \\
        &\Rightarrow \alpha = 0 \\
        \vec{r} \cdot (\vec{c} \times \vec{a}) &\Rightarrow \beta [\vec{a}, \vec{b}, \vec{c}] = 0 \\
        &\Rightarrow \beta = 0 \\
        \vec{r} \cdot (\vec{a} \times \vec{b}) &\Rightarrow \gamma [\vec{a}, \vec{b}, \vec{c}] = 0 \\
        &\Rightarrow \gamma = 0
    \end{align*}
    Hence, $\{\vec{a}, \vec{b}, \vec{c}\}$ form a basis of $\R^3$
\end{proof}

\begin{remark}
    Note that $\vec{a}, \vec{b}, \vec{c}$ need not be unique! As a standard basis in $\R^3$ we use,
    \[
        \{\uvec{i}, \uvec{j}, \uvec{k}\}
    \]
    Where they are
    \begin{enumerate}
        \item \[
            \uvec{i} \cdot \uvec{j} = \uvec{j} \cdot \uvec{k} = \uvec{k} \cdot \uvec{i} = 0  
        \] \hspace*{\fill}(orthogonal)
        \item \[
            \uvec{i} \cdot \uvec{i} = \uvec{j} \cdot \uvec{j} = \uvec{k} \cdot \uvec{k} = 1 
        \] \hspace*{\fill}(unit vectors)
    \end{enumerate}
    Together, we say that these form an \emph{orthonormal} basis.

    Another important point to keep in mind is that the choice of axis is highly arbitrary. Therefore, when we apply vectors to physical laws, we should expect the laws to be invariant with respect to changes of axis. We will cover more about this later in the course when we discuss tensors.
\end{remark}

\subsubsection*{Working in $\R^n$}
\begin{defi}[Linear independance in $\R^n$]
    A set of vectors $\{\vec{v_1}, \vec{v_2}, \cdots, \vec{v_n}\}$ where $v_i \in \R^n$ is \emph{linearly independant} if for scalars $\lambda_i \in \R$
    \begin{equation}\label{eq:2-25}
        \sum_{j=1}^{m}{\lambda_j \vec{v_j}} = \vec{0} \Leftrightarrow \lambda_j = 0 \quad \forall j = 1, 2, \cdots, n
    \end{equation}
\end{defi}

\begin{defi}[Spanning sets in $\R^n$]
    A set of vectors $\{\vec{u_1}, \vec{u_2}, \cdots, \vec{u_m}\}$ where $u_i \in \R^n$ is a \emph{spanning set} if for any $\vec{x} \in \R^n, \quad \exists \lambda_j$ such that,
    \begin{equation}\label{eq:2-26}
        \vec{x} = \sum_{j=1}^{m}{\lambda_j \vec{u_j}}
    \end{equation}
\end{defi}

\begin{remark}
    Note that we cannot span $\R^n$ if $m < n$.
\end{remark}

\begin{defi}[Basis in $\R^n$]
    A linearly independant set of vectors which spans $\R^n$ is a \emph{basis} of $\R^n$. A standard basis that we use
    \[
        \{\vec{e_1}, \vec{e_2}, \cdots, \vec{e_n}\}
    \]
    where
    \[
        e_i = \{00 \cdots 1 \cdots 00\}    
    \]
    where $1$ is in the $i^{th}$ position.
\end{defi}

\begin{defi}[Dimension]
    The \emph{dimension} of a vector space is the number of vectors in the basis.
\end{defi}

\begin{remark}
    We can prove that adding a vector to a basis does not result in a basis, and that if 2 basis exist, then they have the same number of elements. Therefore the dimension of a vector space is well defined.
\end{remark}

We can also define an inner product on this space, for two vectors $\vec{x}, \vec{y} \in \R^n$,
\begin{align*}
    \vec{x} = (x_1, \cdots, x_n) = \sum_{j=1}^{n}{x_i \vec{e_i}} \\
    \vec{y} = (y_1, \cdots, y_n) = \sum_{j=1}^{n}{y_i \vec{e_i}}
\end{align*}
Then the inner product of these two vectors is
\begin{equation}
    \vec{x} \cdot \vec{y} = \braket{x}{y} = \sum_{j=1}^{n}{x_iy_i}
\end{equation}

Note that for this we use the fact that $e_j$ form an orthonormal basis, such that
\[
    e_i \cdot e_j = \begin{cases}
        1 & \text{if} \ i = j \\
        0 & \text{otherwise}
    \end{cases}
\]

\subsubsection*{Working in $\C^n$}

We define the vector space $\C^n$
\[
    \C^n = \{(z_1, z_2, \cdots, z_n) \mid z_i \in \C\}  
\]
as well as a standard basis, which are still
\[
    (1, 0, 0, 0, \cdots), \quad (0, 1, 0, 0, \cdots), \quad (0, 0, 1, 0, \cdots)
\]

\begin{defi}[Scalar product in $\C^n$]
    In $\C^n$ the scalar product is defined, for $\vec{u}, \vec{v} \in \C^n$,
    \[
        \vec{u} \cdot \vec{v} = \braket{u}{v} = \sum_{i=1}^{n}\overline{u_i}v_i \in \C 
    \]
    Note the conjugation of the first argument. We have the following properties
    \begin{enumerate}[eqn]
        \item $\vec{u} \cdot \vec{v} = \overline{\vec{v} \cdot \vec{u}}$ \hspace*{\fill}(Conjugte symmetry)
        \item $\vec{u} \cdot (\lambda \vec{v} + \mu \vec{w}) = \lambda \vec{u} \cdot \vec{v} + \mu \vec{u} \cdot \vec{w}$ \label{eq:2-29}\hspace*{\fill}(Linearity in the $2^{nd}$ argument)
        \item $\vec{u} \cdot \vec{u} = 0 \Leftrightarrow \vec{u} = 0$
    \end{enumerate}

    To solve for the linearity of the first argument we consider
    \begin{align*}
        (\lambda \vec{u} + \mu \vec{v}) \cdot \vec{w} &= \overline{\vec{w} \cdot (\lambda \vec{u} + \mu \vec{v})} \\
        &= \overline{\sqbkt{\lambda \vec{w} \cdot \vec{u} + \mu \vec{w}\vec{v}}} \\
        &= \overline{\lambda} \overline{\vec{w} \cdot \vec{u}} + \overline{\mu} \overline{\vec{w} \cdot \vec{v}} \\
        &= \overline{\lambda} \vec{u} \cdot \vec{w} + \overline{\mu} \vec{v} \cdot \vec{w} \numberthis \label{eq:2-31}
    \end{align*}
\end{defi}

\begin{eg}
    To practice this we will try and expand for $\vec{x}, \vec{y} \in \C^n$   
    \begin{align*}
        &\sum_{k=1}^{4}{(-i)^k \abs{\vec{x} + i^k \vec{y}}} \\
        &= \sum_{k=1}^{4}{(-i)^k \braket{\vec{x} + i^k \vec{y}}{\vec{x} + i^k \vec{y}}} \\
        &= \sum_{k=1}^{4}{(-i)^k \crbkt{\braket{\vec{x} + i^k \vec{y}}{\vec{x}} + i^k\braket{\vec{x} + i^k \vec{y}}{\vec{y}}}} \tag{\ref{eq:2-29}}\\
        &= \sum_{k=1}^{4}{(-i)^k \crbkt{\braket{\vec{x}}{\vec{x}} + \overline{i^k}\braket{\vec{y}}{\vec{x}} + i^k\braket{\vec{x}}{\vec{y}} + i^k \overline{i^k}\braket{\vec{y}}{\vec{y}}}} \tag{\ref{eq:2-31}}\\
        &= \braket{\vec{x}}{\vec{x}}\sum_{k=1}^{4}{(-i)^k}  + \braket{\vec{y}}{\vec{x}}\sum_{k=1}^{4}{(-i)^k\overline{i^k}} + \braket{\vec{x}}{\vec{y}}\sum_{k=1}^{4}{(-i)^ki^k} + \braket{\vec{y}}{\vec{y}}\sum_{k=1}^{4}{(-i)^k i^k \overline{i^k}} \\
        &= \braket{\vec{x}}{\vec{x}}\sum_{k=1}^{4}{(-i)^k}  + \braket{\vec{y}}{\vec{x}}\sum_{k=1}^{4}{(-1)^k} + \braket{\vec{x}}{\vec{y}}\sum_{k=1}^{4}{(1)^k} + \braket{\vec{y}}{\vec{y}}\sum_{k=1}^{4}{(-i)^k}\\
        &= 4\braket{\vec{x}}{\vec{y}}
    \end{align*}
\end{eg}

\subsection{Vector subspaces}
\begin{defi}[Vector subspace]
    A non-empty subset $u$ of vectors of a vector space $V$ is a subspace of $V$ if $u$ is itself a vector space under the same operations as $V$.
\end{defi}

\begin{remark}
    Note that a subspace always includes the additive identity $0$. For instance, a plane is a subspace of 3D space only if it contains the origin.
\end{remark}

\begin{eg}[Trivial examples]
    $V$ and $\{0\}$ are subsapces of $V$. All others are called propper subspaces.
\end{eg}

\begin{thm}
    A subset $u$ of a vector space $V$ is a subspace of $V$ if and only if under the same operations as $V$,
    \begin{enumerate}
        \item $\vec{x}, \vec{y} \in u \Rightarrow \vec{x} + \vec{y} \in u$ \hspace*{\fill}(Closed under addition)
        \item $\vec{x} \in u \Rightarrow \lambda\vec{x}\in u$ for scalar $\lambda$ \hspace*{\fill}(Closed under scalar multiplication)
        \item $\vec{0} \in u$
    \end{enumerate}
    We can condense these three conditions into a single statement. $u$ is a subspace if and only if
    \begin{equation}\label{eq:2-32}
        \lambda \vec{x} + \mu \vec{y} \in u \quad \forall \vec{x}, \vec{y} \in u  
    \end{equation}
    To verify (i) we set $\lambda = \mu = 1$, to verify (ii) we set $\mu = 0$, to verify (iii) we set $\lambda = \mu = 0$
\end{thm}

\begin{eg}\leavevmode
    \begin{enumerate}
    \item Given $\{\vec{a}, \vec{b}, \vec{c}\}$ as a basis of a 3D space, $\{\vec{a} + \vec{c}, \vec{b} + \vec{c}\}$ is a basis of a 2D space.
    Suppose, $\vec{x}, \vec{y} \in \spn\{\vec{a} + \vec{c}, \vec{b} + \vec{c}\}$. Then,
    \begin{align*}
        \vec{x} = \alpha_1 (\vec{a} + \vec{c}) + \beta_1 (\vec{b} + \vec{c}) \\
        \vec{y} = \alpha_2 (\vec{a} + \vec{c}) + \beta_2 (\vec{b} + \vec{c})
    \end{align*}
    Using \eqref{eq:2-32} we get,
    \begin{align*}
        \lambda\vec{x} + \mu \vec{y} =& \lambda\alpha_1 (\vec{a} + \vec{c}) + \lambda\beta_1 (\vec{b} + \vec{c}) \\
        &+ \mu\alpha_2 (\vec{a} + \vec{c}) + \mu\beta_2 (\vec{b} + \vec{c}) \\
        =& (\lambda\alpha_1 + \mu\alpha_2)(\vec{a} + \vec{c}) + (\lambda\beta_1 + \mu\beta_2)(\vec{b} + \vec{c}) \\
        & \in \spn\{\vec{a} + \vec{c}, \vec{b} + \vec{c}\} \quad \therefore \text{subspace}
    \end{align*}
    To show linear independance,
    \begin{align*}
        \alpha (\vec{a} + \vec{c}) + \beta (\vec{b} + \vec{c}) &= 0 \\
        \Rightarrow \alpha\vec{a} + \beta \vec{b} + (\alpha + \beta)\vec{c} &= 0
    \end{align*}
    However since we chose $\{\vec{a}, \vec{b}, \vec{c} \}$ to be a 3D basis, we have linear independance of those vectors,
    \[
        \Rightarrow \alpha = \beta = \alpha + \beta = 0  
    \]
%%%%%%%%%%%%%%%%%%%%%%
% Lecture 6: 18/10/2023
%%%%%%%%%%%%%%%%%%%%%%
    \item $u = \{ \vec{x} \in \R^n \mid \sum_{i=1}^{n}{\alpha_ix_i} = 0, \ \alpha_i \in \R\}$ forms a subspace in $\R^n$.
    Consider $\vec{x}, \vec{y} \in u$.
    \begin{align*}
        \lambda \vec{x} + \mu \vec{y} &= \sum_{i=1}^{n}{\alpha_i(\lambda x_i + \mu y_i)} \\
        &= \lambda\sum_{i=1}^{n}{\alpha_ix_i} + \mu\sum_{i=1}^{n}{y_i} \\
        &= 0
    \end{align*}
    Therefore, $\lambda \vec{x} + \mu \vec{y} \in u$ and $u$ forms a subsapce of $\R^n$. We can think of this subspace as being $n-1$ dimensional. 
    If we think about the restriction imposed by the sum, we are able to freely choose $n-1$ variables before the $n^{th}$ variable is fixed by the sum condition.
    Also if we consider $\alpha = (1, 0, 0, 0, \cdots) = \vec{e_1}$. Then the equation for $u$ is $x_1 = 0$ after which we are free to change all other variables.
    Therefore a basis for this subspace could be $\{\vec{e_2}, \vec{e_3}, \cdots, \vec{e_n}\}$

    \item $w = \{ \vec{x} \in \R^n \mid \sum_{i=1}^{n}{\alpha_ix_i} = 1\}$ is \emph{NOT} a subspace.
    consdier the same linear combination,
    \begin{align*}
        \lambda \vec{x} + \mu \vec{y} &= \sum_{i=1}^{n}{\alpha_i(\lambda x_i + \mu y_i)} \\
        &= \lambda\sum_{i=1}^{n}{\alpha_ix_i} + \mu\sum_{i=1}^{n}{y_i} \\
        &= \lambda + \mu \\
        &\neq 1 \quad \forall \lambda, \mu \in \R
    \end{align*}
    This plane does not pass through the origin, and therefore is not a subspace.
    \end{enumerate}
\end{eg}

\subsection{Einstein Suffix Notation}

Consider a vector $\vec{v} \in \R^3$, we can write
\[
    \vec{v} = v_1\vec{e_1} + v_2\vec{e_2} + v_3\vec{e_3} = (v_1, v_2, v_3)
\]
In such notation, we write the $i^{th}$ component of $\vec{v}$ as $v_i$. Using Einstein notation, we will aim to write out vector equations in terms of these suffixes alone. For instance,
\begin{align*}
    \vec{a} = \vec{c} &\longrightarrow a_i = b_i \\
    \vec{c} = \alpha \vec{a} + \beta \vec{b} &\longrightarrow c_i = \alpha a_i + \beta b_i
\end{align*}
for the scalar product we may write
\[
    \vec{x} \cdot \vec{y} = \sum_{i=1}^{3}{x_iy_i} \longrightarrow x_iy_i
\]
In this case, the presence of a repeated \emph{free} suffix $i$ indicates that we sum over it.

\begin{defi}[Summation Convention]
    When we are writing sums, summation over free suffixes is understood.
    \[
        \vec{x} \cdot \vec{y} = x_iy_i
    \]
    Here $i$ is repeated and called a \emph{dummy} suffix. Since we are summing over it, it matters not what the dummy suffix is called.
\end{defi}

Here are the rules for using Einstein suffix notation.
\begin{enumerate}
    \item If the suffix appears once on the same term: it is a free suffix
    \item If the suffix appears twice in the same term: it is a dummy suffix and we sum over it
    \item If the suffix appears more than twice: this is wrong!
\end{enumerate}

\begin{eg}
    \begin{align*}
        [(\vec{a} \cdot \vec{b})\vec{c} - (\vec{a}\cdot\vec{c}\vec{b})]_i \tag{Note the single free suffix $i$}\\
        =(\vec{a} \cdot \vec{b})\vec{c_i} - (\vec{a}\cdot\vec{c})\vec{b_i} \\
        =a_kb_kc_i - a_kc_kb_i \tag{Note the use of dummy $k$}
    \end{align*}
    Note that for any equation the free suffixes must balance on the left and right. Here we see that we have one free $i$ on the left and right.
\end{eg}

\begin{defi}[Kronecker Delta]
    \[
        \delta_{ij} = \begin{cases}
            1 & \text{if} \ i = j \\
            0 & \text{otherwise}
        \end{cases}
    \]
    It is a symmetric symbol
    \[
        \delta_{ij} = \delta_{ji}
    \]
\end{defi}

\begin{remark}
    Here are some results from using the Kronecker Delta, 
    \begin{align*}
        a_i\delta_{i1} = a_1
    \end{align*}
    In this example we are summing over dummy $i$, however, all terms will be $0$ except for when $i = 1$. In this case the Kronecker delta, pulled out $1$ from the sum. In general,
    \begin{equation}\label{eq:2-33}
        a_i\delta_{ij} = a_j 
    \end{equation}
    We also have
    \begin{equation}\label{eq:2-34}
        \delta_{ij}\delta_{jk} = \delta_{ik}
    \end{equation}
    Where in this example, we have summed the product over all $j$. Note that this sum is only $1$ if $k = i$
    We can also sum over the Kronecker Delta itself,
    \begin{equation}
        \delta_{ii} = 3  
    \end{equation}

    We may compute an example,
    \[
        a_p \delta_{pq} b_q = a \cdot b
    \]
\end{remark}

\begin{defi}[Levi - Cevita (Alternating) Symbol]
    Consider the arrangements of $\{1, 2, 3\}$. We can classify them into two groups
    \[
        \{1, 2, 3\} \quad \{2, 3, 1\} \quad \{3, 1, 2\}
    \]
    where we perform an even number of swaps, and 
    \[
        \{2, 1, 3\} \quad \{1, 3, 2\} \quad \{3, 2, 1\}
    \]
    where we perform an odd number of swaps. These are known as even and odd permutations respectively. The Levi-Cevita symbol is:
    \[
        \varepsilon_{ijk} = \begin{cases}
            +1 & \text{if} \ i, j, k \ \text{is an even permutation of} \ 1, 2, 3 \\
            -1 & \text{if} \ i, j, k \ \text{is an odd permutation of} \ 1, 2, 3 \\
            0 & \text{otherwise}
        \end{cases}
    \]

    We note that $\varepsilon_{ijk}$ has three fixed suffixes.
\end{defi}

\begin{remark}\leavevmode
    \begin{enumerate}
        \item $\varepsilon_{ijk}\delta_{jk} = 0$ since the Kronecker delta is forcing two of the suffixes to be the same, in which case the Levi-Cevita symbol will be $0$.
        \item If $a_{jk} = a_{kj}$ then we say that $a_{jk}$ is symmetric. We can show that
        \[
            \varepsilon_{ijk}a_{jk} = \varepsilon_{ijk}a_{kj} = - \varepsilon_{ikj}a_{kj}
        \]
        Note that we first used the fact that $a_{jk}$ is symmetric, then the fact that when we swap the suffices of $\varepsilon$ once, we move from and even permutation to an odd and vice versa. Lastly, we recall that the names of dummy variables don't matter! Therefore, we get that
        \[
            \varepsilon_{ijk}a_{jk} = 0
        \]
        \item Recall the vector product, 
        \begin{align*}
            \vec{a} \times \vec{b} &= (a_1\uvec{i} + a_2\uvec{j} + a_1\uvec{k}) \times (b_1\uvec{i} + b_2\uvec{j} + b_1\uvec{k}) \\
            &= (a_2b_3 - a_3b_2)\uvec{i} + \cdots \\
            &= \begin{vmatrix}
                \uvec{i} & \uvec{j} & \uvec{k} \\
                a_1 & a_2 & a_3 \\
                b_1 & b_2 & b_3
            \end{vmatrix}
        \end{align*}
        We consider
        \[
            \varepsilon_{1jk}a_jb_k = \sum_{j=1}^{3}\sum_{k=1}^{3} \varepsilon_{1jk}a_jb_k 
        \]
        In this sum, the only non-zero terms are $j=2, k=3$ or $j=3, k=2$. Therefore we get,
        \begin{align*}
            = \varepsilon_{123}a_2b_3 + \varepsilon_{132}a_3b_2 \\
            = a_2b_3 - a_3b_2
        \end{align*}
        In general,
        \begin{equation}
            (\vec{a} \times \vec{b})_i = \varepsilon_{ijk}a_jb_k
        \end{equation}
    \end{enumerate}
\end{remark}

\end{document}